\section{Introdução e Objetivos}
    \subsection{Introdução}
    A cinética química estuda a velocidade das reações químicas e os fatores que a influenciam, como a natureza e a concentração dos reagentes, a temperatura, a superfície de contato e a presença de catalisadores\textsuperscript{[1]}. Segundo a teoria das colisões, uma reação só ocorre quando as partículas dos reagentes colidem com orientação adequada e com energia igual ou superior à energia de ativação ($E_a$). Assim, o aumento da concentração eleva o número de colisões por unidade de tempo, e o aumento da temperatura eleva a energia das partículas e a fração de colisões efetivas, de modo que ambos aumentam a velocidade da reação\textsuperscript{[2]}.

    O relógio de iodo é uma reação adequada para esse estudo porque seu final é indicado de forma repentina pelo aparecimento da cor azul do complexo amido-iodo. O iodato ($\mathrm{IO_3^-}$) reage com o hidrogenossulfito ($\mathrm{HSO_3^-}$) em meio ácido formando iodo ($\mathrm{I_2}$), que é imediatamente consumido enquanto ainda houver hidrogenossulfito na solução. Quando este se esgota, o iodo se acumula e a solução fica azul. Como a quantidade de hidrogenossulfito é a mesma em todos os ensaios, o inverso do tempo até o aparecimento da cor é proporcional à velocidade média da reação\textsuperscript{[1]}.

    \subsection{Objetivos}
    Os objetivos deste experimento foram avaliar o efeito da concentração de iodato e da temperatura sobre a velocidade da reação do relógio de iodo, construir uma curva de calibração do inverso do tempo em função da concentração de $\mathrm{KIO_3}$ e utilizá-la para determinar a concentração de uma solução desconhecida.
